% ========== 2023年03月07号 第01次课 向量 ==========
%\frame{\title{线性代数 \quad 中国科学技术大学 2025春 \\ \ \\ 第一讲 \quad 向量} \date{} \titlepage}




\section*{课程简介}


\begin{frame}[plain]\frametitle{课程考核方式}
\begin{enumerate}
\item 考核成绩为\blue{平时成绩}、\blue{期中成绩}和\blue{期末成绩}加权平均。具体权重会由线性代数课题组根据期中和期末考试的难易度来确定。
\bigskip

\item 平时成绩包含\blue{作业成绩}和\blue{上课考勤}两部分。
\bigskip

\item 周二收发作业。迟交一天到一个星期，得分减半；迟交一个星期以上，零分。（诚信问题零容忍，一旦发现抄袭，将给予零分处理）收发作业务必准时。
\bigskip


\item 习题课为周四第五节课，自愿参加不强制。由助教负责，主要讲解上周布置的习题。
\bigskip

\item 如有事不能来上课，\blue{可以无理由请假, 但必须\underline{课前}告知助教}。
\bigskip

\begin{center}
平时表现 = 5 × 考勤次数 - 请假次数 = \textbf{总评微调力度}。
\end{center}

\end{enumerate}
\end{frame}




\begin{frame}[plain]\frametitle{什么是线性代数(linear algebra)？}
\pp
线性代数是关于\blue{\textbf{向量}空间}和\blue{\textbf{线性}映射}的一个数学分支。
\rightline{------维基百科}
\ \\ \ \\
\pp
线性代数的方法广泛的用在\blue{数学其他分支}、\blue{物理化学}、\blue{计算机科学}、\blue{经济学}等学科中。例如
\begin{enumerate}
\item 泛函分析(研究函数组成的空间);
\item 量子力学(波函数,密度泛函理论);
\item 科学计算(天气预报)；
\item 机器人学(运动学正解)；
\item 数据传输(编码理论)；
\item 人工智能(大模型、向量数据库、深度学习)；
\item $\cdots$
\end{enumerate}\end{frame}




\begin{frame}[plain]
\end{frame}


